\documentclass[12pt]{article}
\usepackage[utf8]{inputenc}
\usepackage[T1]{fontenc}
\usepackage{hyperref}
\title{Maco Language Resources for K-12 Curriculum Development}
\author{Compiled by Claude Sonnet 5.5 (claude-sonnet-5-5) for Andr\'as Kornai}
\date{}
\begin{document}
\maketitle

\section{Introduction}
Maco (also Mako, Wir\"o; self-name Wi\c{}l\"o\c{}; ISO 639-3 wpc; Glottocode maco1239) is a S\'aliban language of Venezuela, spoken mainly in small communities along the middle Ventuari River in Amazonas State. Sources give about 1,000 to 1,200 speakers. The language was almost undocumented, known only from a list of 38 words collected around 1900, until Jorge Ros\'es Labrada's 2015 work. The name ``Mako'' has also been applied to other, distinct groups in the area, so take care not to confuse them.

\section{Primary Educational Resources}
No primer, textbook or school programme for Maco appeared in search results.

\section{Linguistic and Reference Materials}
\begin{itemize}
\item Jorge Ros\'es Labrada, \emph{The Mako language: Vitality, Grammar and Classification} (PhD dissertation, University of Western Ontario, 2015): \url{https://ir.lib.uwo.ca/etd/2851}
\item Record of the same dissertation at CORE: \url{https://core.ac.uk/works/9175237}
\item Collaborative Documentation and Description of Mako, a S\'aliban Language of Venezuela (ELDP archive collection): \url{https://eldp.access.preservica.com/dk0174/}
\item Grambank entry for Maco: \url{https://grambank.clld.org/languages/maco1239}
\item Wikipedia article on the language (as Wir\"o): \url{https://en.wikipedia.org/wiki/Wir%C3%B6_language}
\end{itemize}
No dictionary was found.

\section{Religious and Community Resources}
A Scripture Earth page exists for code wpc, but the search result showed only its title, so it is not known what, if anything, it holds: \url{https://scriptureearth.org/00eng.php?iso=wpc}

\section{Digital and Audio Resources}
The ELDP collection above is the documentation archive. It reportedly includes recordings made in the Ventuari River communities. Its contents were not examined.

\section{Curriculum Development Context}
The language is still spoken, but no school materials were found. The 2015 grammar and the ELDP collection are the starting points for developing a primer and orthography. Any such work should be done with Mako communities and with the documenting linguist.

\end{document}